\documentclass[11pt,a4paper]{article}

\usepackage[utf8]{inputenc}
\usepackage[T1]{fontenc}
\usepackage[spanish,provide=*]{babel}
\usepackage{csquotes}
\usepackage{charter}
\usepackage[scaled=0.92]{helvet}
\usepackage[a4paper,top=2.6cm,bottom=2.6cm,left=2.4cm,right=2.4cm,headheight=14pt]{geometry}
\usepackage{booktabs}
\usepackage{longtable}
\usepackage{array}
\usepackage{xcolor}
\usepackage{colortbl}
\usepackage{titlesec}
\usepackage{fancyhdr}
\usepackage{enumitem}
\usepackage{tcolorbox}
\usepackage{listings}
\usepackage{amssymb}
\usepackage{microtype}
\usepackage[hidelinks]{hyperref}

\definecolor{azulUTEQ}{HTML}{123A5F}
\definecolor{azulClaro}{HTML}{E8EEF4}
\definecolor{grisTexto}{HTML}{3A3A3A}
\definecolor{acento}{HTML}{A6242B}
\definecolor{grisCodigo}{HTML}{F2F4F6}

\hypersetup{
	colorlinks=true,
	linkcolor=azulUTEQ,
	urlcolor=azulUTEQ,
	citecolor=azulUTEQ
}

\titleformat{\section}{\normalfont\Large\bfseries\color{azulUTEQ}}{\thesection}{0.7em}{}
\titleformat{\subsection}{\normalfont\large\bfseries\color{azulUTEQ}}{\thesubsection}{0.6em}{}
\titlespacing*{\section}{0pt}{18pt}{8pt}
\titlespacing*{\subsection}{0pt}{14pt}{6pt}

\pagestyle{fancy}
\fancyhf{}
\renewcommand{\headrulewidth}{0.4pt}
\fancyhead[L]{\footnotesize\color{azulUTEQ}TA-PFC-E4 $\cdot$ LabQuorum (equipo FUVV) $\cdot$ Aplicaciones Distribuidas (ISR-701)}
\fancyhead[R]{\footnotesize\color{azulUTEQ}\thepage}

\setlength{\parskip}{6pt}
\setlength{\parindent}{0pt}
\renewcommand{\arraystretch}{1.18}

\lstset{
	basicstyle=\ttfamily\footnotesize,
	backgroundcolor=\color{grisCodigo},
	frame=single,
	rulecolor=\color{azulUTEQ},
	breaklines=true,
	columns=fullflexible,
	xleftmargin=6pt,
	xrightmargin=6pt,
	framexleftmargin=6pt
}

\newtcolorbox{aviso}[1]{
	colback=azulClaro,
	colframe=azulUTEQ,
	boxrule=0.8pt,
	arc=2pt,
	left=8pt,
	right=8pt,
	top=7pt,
	bottom=7pt,
	title={\bfseries #1},
	fonttitle=\normalsize\color{white},
	colbacktitle=azulUTEQ
}

\newtcolorbox{alerta}[1]{
	colback=white,
	colframe=acento,
	boxrule=0.9pt,
	arc=2pt,
	left=8pt,
	right=8pt,
	top=7pt,
	bottom=7pt,
	title={\bfseries #1},
	fonttitle=\normalsize\color{white},
	colbacktitle=acento
}

\begin{document}

\begin{titlepage}
	\centering
	\vspace*{1.2cm}
	{\color{azulUTEQ}\rule{\textwidth}{2pt}}\\[10pt]
	{\LARGE\bfseries\color{azulUTEQ} Guía y rúbrica de la entrega final TA-PFC-E4\par}
	\vspace{8pt}
	{\Large\bfseries\color{azulUTEQ} PFC C --- LabQuorum (equipo FUVV)\par}
	\vspace{4pt}
	{\normalsize\color{grisTexto}\emph{denominación anterior: Control y Administración de Laboratorios Informáticos}\par}
	\vspace{10pt}
	{\color{azulUTEQ}\rule{\textwidth}{2pt}}\\[16pt]
	{\large Asignatura: Aplicaciones Distribuidas (ISR-701)\par}
	\vspace{4pt}
	{\large Carrera de Ingeniería de Software $\cdot$ Universidad Técnica Estatal de Quevedo\par}
	\vspace{24pt}
	{\large\bfseries Entrega acumulativa total y paquete de evidencias\par}
	{\large\bfseries orientado a publicación en revista JCR híbrida\par}
	\vspace{24pt}
	\begin{minipage}{0.86\textwidth}
		\centering
		\color{grisTexto}
		\small
		\textbf{Revista objetivo:} Future Generation Computer Systems (Elsevier), ISSN 0167-739X.\\
		JIF 2025 = 5,9 $\cdot$ Q1 en \emph{Computer Science, Theory \& Methods} $\cdot$ SCIE $\cdot$ modelo híbrido.\\[6pt]
		Eje del alcance editorial: \emph{Applications and Application Support}, que nombra expresamente los laboratorios virtuales y los sistemas colaborativos.\\[6pt]
		Datos bibliométricos verificados en Journal Citation Reports el 25 de agosto de 2026; conjunto de datos actualizado el 17 de junio de 2026.
	\end{minipage}
	\vfill
	{\small Quevedo, Ecuador $\cdot$ 25 de agosto de 2026\par}
\end{titlepage}

\tableofcontents
\newpage

\section{Regla de esta entrega}

TA-PFC-E4 es una \textbf{entrega total}: el sistema completo, la documentación completa y el paquete completo de evidencias, con independencia de lo entregado en E1, E2 o E3. Si un artefacto no está en este paquete, no existe a efectos de calificación.

\begin{aviso}{Este es el proyecto con mejor posición de partida}
	De los cuatro PFC de la asignatura, este es el único cuya guía ya le atribuye una pregunta de investigación propia: cómo mejorar el acceso concurrente a los recursos del laboratorio. Tiene un fenómeno medible, un mecanismo que proponer, alternativas contra las que compararlo y una prueba de carga que lo valida o lo refuta. Por eso su revista objetivo es la más exigente de las cuatro que se han asignado, y por eso la exigencia experimental de esta guía es también la más alta.
\end{aviso}

\section{Del nombre a la publicabilidad}

\subsection{Por qué el nombre actual no sirve}

«Control y Administración de Laboratorios Informáticos» describe una función administrativa, no una contribución. Nombra el dominio y no el mecanismo; es una frase, no un nombre; no se puede usar como identificador de repositorio, de paquete ni de organización; y traducido al inglés produce un título tan genérico que un editor lo asocia de inmediato a un trabajo de gestión, no de sistemas distribuidos.

\subsection{Un ejemplo real de la comprobación que hay que hacer}

Al buscar nombre para este proyecto se descartaron dos candidatos por colisión, y conviene que el equipo sepa por qué, porque tendrá que repetir esta comprobación:

\begin{itemize}[leftmargin=*,itemsep=3pt]
	\item \textsc{SlotGuard} se descartó: ya existe un producto comercial con ese nombre, listado en directorios de software.
	\item \textsc{LabLease} se descartó: existe un repositorio público con ese nombre exacto.
\end{itemize}

La comprobación es obligatoria y se documenta: buscador general, GitHub y directorios de software. Un nombre colisionado obliga a renombrar después de publicar, que es cuando ya no se puede.

\subsection{Nombre adoptado}

\begin{aviso}{Denominación oficial a partir de esta entrega}
	\textbf{Nombre del sistema:} \textsc{LabQuorum}\\[6pt]
	\textbf{Descriptor:} arbitraje de reservas concurrentes sobre recursos compartidos de laboratorio.\\[6pt]
	\textbf{Título previsto del manuscrito:} \emph{LabQuorum: arbitration of concurrent reservations for shared laboratory resources under failure --- design and empirical evaluation}.
\end{aviso}

El nombre no es decorativo: alude al mecanismo de coordinación que el sistema debe implementar y evaluar, y por eso el equipo está obligado a que una de las estrategias comparadas sea efectivamente de consenso por quórum. Un nombre que promete algo que el sistema no hace es un problema de honestidad, no de mercadotecnia. Alternativas admisibles si se comprueba que están libres: \textsc{LabArbiter} o \textsc{CronoLab}.

\section{La revista objetivo y por qué es esta}

\subsection{Identificación}

\begin{center}
	\small
	\begin{tabular}{@{}p{4.4cm}p{10.6cm}@{}}
		\toprule
		\textbf{Revista} & Future Generation Computer Systems (FGCS), Elsevier \\
		\textbf{ISSN} & 0167-739X (impreso) $\cdot$ 1872-7115 (electrónico) \\
		\textbf{Indexación} & Web of Science, SCIE \\
		\textbf{JIF 2025} & 5,9 \\
		\textbf{Cuartil JCR} & Q1 en \emph{Computer Science, Theory \& Methods} \\
		\textbf{Modelo} & Híbrido: la vía de suscripción no tiene coste para el autor \\
		\bottomrule
	\end{tabular}
\end{center}

\subsection{Las cuatro razones de la elección}

\begin{enumerate}[leftmargin=*,itemsep=5pt]
	\item \textbf{Es híbrida y es Q1}, con el factor de impacto más alto de las cuatro revistas objetivo asignadas en la asignatura.
	\item \textbf{Su alcance nombra expresamente los laboratorios.} El eje \emph{Applications and Application Support} de FGCS incluye los laboratorios virtuales y los sistemas colaborativos entre sus temas. No hay que forzar el encaje: el dominio del proyecto está escrito en el alcance de la revista.
	\item \textbf{Publica exactamente este tipo de mecanismo.} Hay trabajo reciente en la propia revista sobre orquestación distribuida basada en consenso para microservicios en clústeres de borde \cite{morabito2026consensus} y sobre protocolos de control de concurrencia optimista \cite{wang2024dodo}. El estado del arte del manuscrito se ancla en la propia revista, que es lo que un editor espera.
	\item \textbf{Su eje declarado es aplicado.} FGCS busca aplicaciones novedosas sobre infraestructuras distribuidas, no solo teoría. Un sistema real, desplegado y medido es el perfil de artículo que publica.
\end{enumerate}

\begin{alerta}{Honestidad sobre las probabilidades}
	FGCS recibe un volumen enorme de manuscritos y su tasa de rechazo de escritorio es alta. Que este PFC sea el mejor posicionado de los cuatro no lo convierte en una apuesta segura. La escalera de respaldo no es un plan de consolación: es parte del plan.
\end{alerta}

\subsection{Escalera de respaldo}

\begin{center}
	\small
	\begin{tabular}{@{}p{5.4cm}rp{1.5cm}p{2.0cm}p{4.3cm}@{}}
		\toprule
		\textbf{Revista} & \textbf{JIF} & \textbf{Cuartil} & \textbf{Editorial} & \textbf{Cuándo redirigir} \\
		\midrule
		Journal of Network and Computer Applications & 7,7 & Q1 & Elsevier & Acepta también notas breves: útil si la contribución resulta acotada \\
		\addlinespace
		Journal of Systems Architecture & 5,3 & Q1 & Elsevier & Si el peso recae en los agentes por laboratorio y la arquitectura del sistema \\
		\addlinespace
		Concurrency and Computation: Practice and Experience & 2,2 & Q3 & Wiley & La más receptiva a trabajos aplicados de sistemas distribuidos; recurso final \\
		\bottomrule
	\end{tabular}
\end{center}

Las tres son híbridas y están verificadas en JCR. Las evidencias de esta guía son la unión de lo que piden las cuatro.

\section{La contribución publicable de LabQuorum}

\subsection{El fenómeno}

Dos personas piden el mismo equipo en la misma franja horaria en el mismo instante. El sistema tiene que otorgárselo a una sola. Eso parece trivial hasta que el sistema está distribuido entre varios laboratorios, con agentes que reportan estado, y hasta que el nodo que arbitra se cae a mitad de la ráfaga. Entonces el problema deja de ser de programación y pasa a ser de coordinación distribuida, que es lo que la asignatura enseña y lo que una revista publica.

\subsection{La contribución}

LabQuorum no aporta un algoritmo nuevo. Aporta algo que sí falta y que una revista aplicada valora: la \textbf{comparación empírica, en un sistema real desplegado, de cinco estrategias de arbitraje de reservas concurrentes sobre un recurso físico compartido, evaluadas conjuntamente en seguridad, latencia, equidad y comportamiento ante la caída del árbitro}. El compromiso entre consistencia y disponibilidad que subyace está formalizado \cite{gilbert2002brewer}; la relación entre confirmación atómica y consenso, también \cite{gray2006consensus}; y la ordenación causal sin reloj compartido es un resultado clásico \cite{lamport1978time}. Lo que no está medido es cuánto cuesta cada opción en este escenario concreto.

\subsection{Preguntas de investigación}

\begin{aviso}{Las tres preguntas de LabQuorum}
	\textbf{PI-1 (seguridad, la principal).} ¿Qué estrategias de arbitraje garantizan cero dobles adjudicaciones bajo contención máxima sobre la misma franja, y cuáles fallan y con qué frecuencia?\\[6pt]
	\textbf{PI-2 (coste).} ¿Qué precio paga cada estrategia que sí garantiza la seguridad, en latencia de adjudicación, rendimiento y rechazos innecesarios, al crecer la contención?\\[6pt]
	\textbf{PI-3 (resiliencia).} ¿Qué ocurre cuando el nodo coordinador cae a mitad de una ráfaga de solicitudes: cuánto tarda la reelección, cuántas solicitudes se pierden o se duplican, y se mantiene la garantía de seguridad durante esa ventana?
\end{aviso}

PI-1 es la pregunta de investigación que la guía de la asignatura ya le atribuía al proyecto, ahora formulada de modo que admite una respuesta negativa. PI-3 es la que eleva el trabajo por encima de un ejercicio de aula: casi nadie mide qué pasa durante la reelección.

\section{Diseño experimental exigido}

\subsection{Las cinco estrategias comparadas}

\begin{center}
	\small
	\begin{tabular}{@{}p{1.1cm}p{4.6cm}p{9.3cm}@{}}
		\toprule
		\textbf{Cód.} & \textbf{Estrategia} & \textbf{Qué implementa} \\
		\midrule
		S0 & Sin arbitraje & Línea base: escritura directa. Debe producir dobles adjudicaciones y sirve para demostrar que el fenómeno existe \\
		\addlinespace
		S1 & Bloqueo optimista & Control por versión: se detecta el conflicto al confirmar y se rechaza al perdedor \\
		\addlinespace
		S2 & Bloqueo pesimista & Se toma el candado sobre la franja antes de decidir \\
		\addlinespace
		S3 & Coordinador elegido + orden total & Un nodo líder, elegido por el algoritmo del matón (\emph{Bully}), serializa las solicitudes usando marcas lógicas de Lamport \cite{lamport1978time} \\
		\addlinespace
		S4 & Transacción serializable por consenso & La adjudicación se decide dentro de una transacción serializable respaldada por consenso de quórum en el clúster de datos \\
		\bottomrule
	\end{tabular}
\end{center}

Las cinco deben ser conmutables mediante una variable de entorno. No se admiten cinco ramas divergentes: eso impide comparar y hace la medición irreproducible.

\subsection{Experimento 1 --- Seguridad y coste bajo contención}

Diseño factorial completo: 5 estrategias $\times$ 3 niveles de contención (10, 50 y 200 solicitudes simultáneas \emph{sobre la misma franja del mismo equipo}) $\times$ 2 escenarios (nominal y con caída del coordinador a mitad de la ráfaga) $=$ 30 condiciones. Cinco repeticiones independientes por condición, con reinicio completo del entorno. Total: \textbf{150 ejecuciones}. Como cada ráfaga dura segundos, el experimento completo cabe en una jornada.

\begin{center}
	\small
	\begin{tabular}{@{}p{5.4cm}p{3.2cm}p{6.4cm}@{}}
		\toprule
		\textbf{Variable de respuesta} & \textbf{Unidad} & \textbf{Por qué se mide} \\
		\midrule
		\textbf{Tasa de doble adjudicación} & \% de franjas & \textbf{Es el resultado central: la violación de seguridad} \\
		Latencia de adjudicación p50, p95, p99 & milisegundos & Precio que paga el usuario por la garantía \\
		Rendimiento & adjudicaciones por segundo & Capacidad efectiva bajo contención \\
		Rechazos innecesarios & \% de solicitudes & Penaliza al optimista: rechaza sin que hubiera conflicto real \\
		Índice de equidad de Jain & adimensional & Detecta si una estrategia favorece sistemáticamente a un nodo o a un laboratorio \\
		Consumo de CPU y memoria & \% y MiB & Coste de infraestructura de cada estrategia \\
		\bottomrule
	\end{tabular}
\end{center}

\subsection{Experimento 2 --- Recuperación del árbitro}

Solo aplica a S3 y S4. Se provoca la caída del coordinador en tres momentos distintos de la ráfaga (al 25\,\%, al 50\,\% y al 75\,\% de las solicitudes), cinco repeticiones cada uno. Se mide: tiempo de detección de la caída, tiempo de reelección, número de solicitudes perdidas, número de solicitudes duplicadas, y si se mantuvo la garantía de cero dobles adjudicaciones durante la ventana.

\subsection{Experimento 3 --- Monitoreo bajo el modelo de fallos real}

El punto débil documentado del proyecto son los reportes de estado de los agentes de laboratorio, que fallan por omisión y por temporización. Se mide, con 20 agentes simulados y tres tasas de pérdida de reportes (0\,\%, 5\,\% y 20\,\%): tiempo de detección de un equipo caído mediante latidos, tasa de falsos positivos de indisponibilidad, y desfase entre el estado real y el estado publicado. Este experimento sostiene la afirmación de que el módulo de monitoreo puede posicionarse como AP mientras el de reservas es CP.

\subsection{Los invariantes que define el oráculo}

Un \emph{oráculo de consistencia} es un programa que, al terminar cada ejecución, revisa la base de datos y responde sí o no a la pregunta «¿quedó todo cuadrado?». Sin él no se puede medir la variable central. LabQuorum debe verificar al menos:

\begin{enumerate}[leftmargin=*,itemsep=3pt]
	\item Ningún equipo tiene dos reservas confirmadas que se solapen en el tiempo.
	\item Toda reserva confirmada tiene exactamente un usuario adjudicatario.
	\item Ninguna franja fue adjudicada sobre un equipo marcado en mantenimiento.
	\item Toda reserva cancelada liberó la franja exactamente una vez, ni cero ni dos.
	\item El registro de control de acceso no permite entrada sin reserva vigente.
\end{enumerate}

\subsection{Análisis estadístico}

Para toda comparación se exige mediana con intervalo de confianza del 95\,\%, prueba U de Mann-Whitney (compara si un grupo tiende a dar valores mayores que otro sin suponer distribución normal), tamaño del efecto mediante el estadístico $\hat{A}_{12}$ de Vargha-Delaney, y diagramas de caja por condición. Las tasas de doble adjudicación se reportan con intervalo de confianza binomial: un cero observado en 50 ejecuciones no es lo mismo que un cero garantizado, y el manuscrito debe decirlo con esas palabras.

\subsection{Amenazas a la validez}

\begin{center}
	\small
	\begin{tabular}{@{}p{3.2cm}p{11.8cm}@{}}
		\toprule
		\textbf{Tipo} & \textbf{Qué debe discutir el equipo} \\
		\midrule
		Interna & Ruido de la máquina, interferencia entre contenedores, sincronía artificial de las ráfagas, semillas de la inyección de fallos \\
		\addlinespace
		Externa & Número de laboratorios y equipos simulados frente a un campus real; patrón de demanda sintético frente al real \\
		\addlinespace
		De constructo & Si «doble adjudicación» capta todo lo que importa, y si el índice de Jain es la medida de equidad adecuada para este dominio \\
		\addlinespace
		De conclusión & Repeticiones, potencia estadística y comparaciones múltiples entre cinco estrategias \\
		\bottomrule
	\end{tabular}
\end{center}

\section{Evidencias obligatorias}

\begin{center}
	\footnotesize
	\begin{longtable}{@{}p{0.9cm}p{4.6cm}p{4.2cm}p{5.0cm}@{}}
		\toprule
		\textbf{Cód.} & \textbf{Evidencia} & \textbf{Artefacto entregable} & \textbf{Criterio de aceptación} \\
		\midrule
		\endhead
		E01 & Repositorio público renombrado & URL de GitHub & Nombre \textsc{LabQuorum}; los cuatro integrantes con \emph{commits} sustantivos en al menos seis semanas distintas \\
		E02 & Comprobación del nombre & Nota en el \texttt{README.md} & Evidencia de la búsqueda de colisiones en buscador, GitHub y directorios \\
		E03 & Despliegue reproducible & \texttt{docker-compose.yml} & El sistema completo, incluidos los agentes simulados, levanta con un solo comando \\
		E04 & Estrategia conmutable & \texttt{ARBITER=s0$\mid$s1$\mid$s2$\mid$s3$\mid$s4} & Cambiar el valor y reiniciar basta para conmutar entre las cinco \\
		E05 & Generador de ráfagas & Script de carga determinista & Produce N solicitudes verdaderamente simultáneas sobre la misma franja; misma semilla, misma secuencia \\
		E06 & Inyector de fallos & Herramienta de caída del coordinador y de pérdida de reportes & Momento de la caída y tasa de pérdida configurables con semilla \\
		E07 & Simulador de agentes & 20 agentes de laboratorio & Reportan estado por TCP con latidos; tasa de pérdida parametrizable \\
		E08 & Oráculo de consistencia & Script verificador & Comprueba los cinco invariantes y emite informe por ejecución \\
		E09 & Datos crudos & CSV de las 150 ejecuciones y de los experimentos 2 y 3 & Las figuras del documento se regeneran ejecutando el \emph{notebook} \\
		E10 & Tabla de seguridad & Matriz estrategia $\times$ contención & Tasa de doble adjudicación con intervalo de confianza binomial \\
		E11 & Medición de recuperación & Tabla de reelección & Tiempo de detección, de reelección, solicitudes perdidas y duplicadas \\
		E12 & Cómputo paralelo & Job PySpark + medición & Análisis de los registros de uso de laboratorios; \emph{speedup} con 1, 2, 4 y 8 \emph{workers} y fracción de Amdahl calculada \\
		E13 & Trazas distribuidas & Captura exportada & Recorrido completo de una solicitud de reserva en contención \\
		E14 & Observabilidad & Tablero Grafana + métricas & Estado en vivo de todos los laboratorios; versionado como código \\
		E15 & Pruebas & Informe de cobertura & Cobertura $\geq 70\,\%$ del motor de reservas; integración y carga \\
		E16 & Calidad del producto & Evaluación ISO/IEC 25010:2023 & Al menos cuatro características con métrica real, priorizando fiabilidad y eficiencia de desempeño \cite{iso25010_2023} \\
		E17 & Documento del PFC & \texttt{.tex} + PDF & Compila desde el repositorio clonado; bibliografía IEEE con DOI que resuelven \\
		E18 & Borrador del manuscrito & \texttt{.tex} + PDF & Título y estructura de la sección 9; máximo 8.000 palabras \\
		E19 & Declaración de uso de IA & Sección del documento & Herramienta, propósito, alcance y qué revisó cada integrante \\
		E20 & Paquete de replicación & DOI de Zenodo & El DOI resuelve; contiene E05 a E09 \\
		E21 & Carátula de entrega & PDF subido al SGA & Datos de identificación y la URL del repositorio en una sola línea \\
		\bottomrule
	\end{longtable}
\end{center}

\section{Cobertura de los temas de la asignatura}

\begin{center}
	\footnotesize
	\begin{longtable}{@{}p{1.2cm}p{5.2cm}p{8.4cm}@{}}
		\toprule
		\textbf{Unidad} & \textbf{Tema} & \textbf{Evidencia concreta exigida en LabQuorum} \\
		\midrule
		\endhead
		U1 & Transparencias ANSA & Al menos cinco de las ocho, cada una con el mecanismo que la realiza \\
		U1 & Sockets TCP & Evidencia E07: agentes que reportan estado y latidos con delimitación por longitud \\
		U1 & gRPC / RPC & Contratos \texttt{.proto} entre Reservas, Monitoreo y Control-Acceso \\
		U1 & Relojes de Lamport y vectoriales & \textbf{Núcleo del artículo}: estrategia S3, evidencias E04, E08 y E10 \cite{lamport1978time} \\
		U1 & Tolerancia a fallos & Evidencias E06 y E11; modelo de fallos por omisión y temporización de los agentes \\
		U1 & Elección de líder & \textbf{Medida, no solo implementada}: evidencia E11 con tiempos de detección y reelección \\
		U1 & Teorema CAP & Posición por agregado: CP en reservas, AP en monitoreo, sostenida con la evidencia del experimento 3 \cite{gilbert2002brewer} \\
		U1 & Coordinación y consistencia & Estrategias S2, S3 y S4; consenso de quórum en el clúster \cite{gray2006consensus} \\
		U1 & Seguridad & JWT, control de acceso por cuatro roles, TLS y enlace entre acceso físico y reserva vigente \\
		U1 & SLA y alta disponibilidad & Objetivo declarado en horas pico y su cumplimiento medido \\
		\addlinespace
		U3 & Microservicios con base propia & Diagrama C4 nivel 3 y justificación de fronteras \\
		U3 & API REST + OpenAPI & Especificación OpenAPI 3.0 validada y versionada \\
		U3 & API Gateway & Enrutamiento, autenticación y limitación de tasa \\
		U3 & Mensajería asíncrona & Eventos \texttt{EquipoReservado} y \texttt{EquipoEnFalla}; justificación del \emph{broker} \\
		U3 & Docker y orquestación & Evidencia E03; manifiestos de Kubernetes para Reservas y Monitoreo \cite{deslauriers2025deployment} \\
		U3 & Patrones de despliegue & Demostración de una actualización sin interrumpir el servicio de reservas \\
		\addlinespace
		U2 & Fragmentación & Particionado por laboratorio, edificio o campus \\
		U2 & Replicación y consistencia & Clúster CockroachDB con consistencia serializable; es la base de S4 \\
		U2 & Tolerancia a fallos en datos & Prueba documentada de caída de un nodo sin perder reservas confirmadas \\
		\addlinespace
		U4 & Procesamiento distribuido & Evidencia E12: ocupación por hora y equipo, software más usado, equipos con fallas recurrentes \\
		U4 & Ley de Amdahl & Fracción secuencial estimada de la curva medida, no supuesta \\
		\addlinespace
		U5 & N-capas y SOLID & Estructura de capas por microservicio con responsabilidades únicas \\
		U5 & Patrones GoF & State para el estado del equipo, Observer para el monitoreo, Strategy para las políticas de reserva, Repository y Proxy \\
		U5 & Inyección de dependencias & Contenedor IoC con ejemplo de sustitución en pruebas \\
		U5 & CI/CD & Flujo de GitHub Actions con pruebas, imagen y análisis de calidad \\
		U5 & Observabilidad & Evidencias E13 y E14 \\
		U5 & Pruebas & Evidencia E15 \\
		U5 & Calidad ISO/IEC 25010 & Evidencia E16 \\
		\bottomrule
	\end{longtable}
\end{center}

\section{Del documento del PFC al manuscrito}

\begin{center}
	\small
	\begin{tabular}{@{}p{3.8cm}p{1.7cm}p{9.3cm}@{}}
		\toprule
		\textbf{Sección} & \textbf{Extensión} & \textbf{De dónde sale y qué debe contener} \\
		\midrule
		Introducción & 1,5 pág. & El problema de la adjudicación concurrente de recursos físicos compartidos. Termina enunciando PI-1, PI-2 y PI-3 \\
		\addlinespace
		Trabajo relacionado & 1,5 pág. & De 15 a 25 referencias, con al menos cinco de la propia revista objetivo. Anclarlo a reservas, laboratorios y coordinación distribuida, nunca a control industrial ni a redes eléctricas inteligentes \\
		\addlinespace
		Diseño de LabQuorum & 2 pág. & Arquitectura, agentes por laboratorio y las cinco estrategias de arbitraje \\
		\addlinespace
		Diseño del estudio & 2,5 pág. & Sección 5 de esta guía: los tres experimentos, invariantes y análisis \\
		\addlinespace
		Resultados & 3 pág. & De E09, E10 y E11. Solo hechos medidos \\
		\addlinespace
		Discusión y lecciones aprendidas & 2 pág. & Qué estrategia recomendarían a una universidad real y a qué coste \\
		\addlinespace
		Amenazas a la validez & 0,7 pág. & Sección 5.7 de esta guía \\
		\addlinespace
		Conclusiones & 0,5 pág. & Respuesta explícita a las tres preguntas \\
		\addlinespace
		Disponibilidad de datos & 3 líneas & DOI de Zenodo de E20 \\
		\bottomrule
	\end{tabular}
\end{center}

\begin{aviso}{La advertencia sobre las referencias que ya venía en la guía de la asignatura}
	El estado del arte debe anclarse al dominio: gestión académica, reserva de recursos, laboratorios y sistemas educativos distribuidos. Citar trabajos de control industrial o de redes eléctricas inteligentes porque «también son sistemas distribuidos» es la señal más rápida de que el equipo no leyó lo que cita, y un revisor la detecta en la primera página.
\end{aviso}

\section{Cronograma de las semanas finales}

\begin{center}
	\small
	\begin{tabular}{@{}p{1.8cm}p{6.4cm}p{6.6cm}@{}}
		\toprule
		\textbf{Semana} & \textbf{Objetivo} & \textbf{Hito verificable} \\
		\midrule
		14 & Renombrado a \textsc{LabQuorum} & Repositorio, código, contenedores y documento renombrados; comprobación de colisiones documentada \\
		15 & Cerrar E04, E05 y E06 & Las cinco estrategias conmutan; el generador de ráfagas y el inyector funcionan \\
		15 & Cerrar E07 y E08 & Agentes simulados y oráculo operativos; ejecución piloto completa \\
		16 & Experimentos 1, 2 y 3 & Datos crudos de E09, tabla de seguridad E10 y de recuperación E11 \\
		16 & Cerrar E12 a E16 & Spark, trazas, tablero, pruebas y calidad \\
		17 & Cerrar E17 y E19 & Documento compilando desde clonación limpia \\
		17 & Cerrar E18 y E20 & Borrador del manuscrito y DOI de Zenodo \\
		17 & Defensa oral & Demostración en vivo de la contención y de la caída del coordinador \\
		\bottomrule
	\end{tabular}
\end{center}

\begin{alerta}{Riesgo de planificación}
	El experimento 2 exige provocar la caída del coordinador en un momento controlado de la ráfaga. Es la pieza técnicamente más delicada de todo el proyecto y la que con más probabilidad se subestima. Debe estar funcionando al final de la semana 15, no al principio de la 16.
\end{alerta}

\section{Preguntas de la defensa oral}

\begin{enumerate}[leftmargin=*,itemsep=4pt]
	\item ¿Cómo resuelven, con relojes lógicos, dos reservas concurrentes sobre el mismo equipo y la misma franja?
	\item ¿Cuál de las cinco estrategias produjo dobles adjudicaciones y en qué proporción? Muestren el dato.
	\item ¿Qué precio en latencia paga la estrategia más segura frente a la más rápida, y cuál es el tamaño del efecto?
	\item ¿Por qué el módulo de reservas se posiciona como CP y el de monitoreo puede ser AP? ¿Qué evidencia propia lo respalda?
	\item Si el coordinador cae con cien solicitudes en vuelo, ¿cuántas se pierden, cuántas se duplican y en cuánto tiempo se recupera el servicio?
	\item ¿Qué mide su índice de equidad y por qué importa en un laboratorio universitario?
	\item ¿Qué fracción del análisis de registros resultó paralelizable y cómo lo dedujeron de la curva medida?
	\item Si tuvieran que desplegar LabQuorum en la UTEQ el próximo semestre, ¿qué estrategia elegirían y con qué dato lo justifican?
\end{enumerate}

\section{Rúbrica de evaluación}

\subsection{Criterios de piso}

\begin{alerta}{Incumplir cualquiera de estos criterios implica calificación CERO en toda la entrega}
	\textbf{P1.} El estudiante debe subir al SGA un documento PDF con carátula que muestre claramente, en una sola línea, la URL del repositorio donde se encuentra el trabajo desarrollado. Si la URL está rota, el repositorio no es accesible o no se cumple esta directriz, la calificación es CERO.\\[6pt]
	\textbf{P2.} Si el PDF no se genera correctamente a partir del LaTeX mediante clonación del repositorio y compilación del archivo \texttt{.tex}, la calificación es CERO. Las instrucciones de compilación (compilador, orden de comandos, archivo principal, dependencias) deben estar documentadas en el archivo \texttt{README.md} del repositorio.\\[6pt]
	\textbf{P3.} Si el sistema no levanta con un solo comando desde una máquina limpia siguiendo el \texttt{README.md}, la calificación es CERO.\\[6pt]
	\textbf{P4.} Si no existen los datos crudos de E09 o el oráculo de E08, la calificación es CERO: sin ellos la entrega no tiene resultado medido.\\[6pt]
	\textbf{P5.} Si la estrategia S0 no produce ninguna doble adjudicación bajo contención máxima, la calificación es CERO por experimento mal construido: significa que las solicitudes no fueron realmente simultáneas y, por tanto, ninguna medición del resto de estrategias es válida.\\[6pt]
	\textbf{P6.} Si se detectan referencias fabricadas, DOI que no resuelven o citas cuyos campos no corresponden a la obra real, la calificación es CERO por falta de integridad académica.\\[6pt]
	\textbf{P7.} Si falta la declaración de uso de IA (E19), la calificación es CERO.
\end{alerta}

\begin{aviso}{Por qué P5 es un criterio de piso y no una observación}
	Es el control negativo del experimento. Si la línea base sin arbitraje no falla nunca, no se está midiendo contención: se están enviando solicitudes que en realidad llegan una detrás de otra. Un revisor lo detecta de inmediato y, con él, se cae todo el artículo. Comprobarlo es lo primero que debe hacer el equipo tras la ejecución piloto.
\end{aviso}

\subsection{Criterios calificables}

\begin{center}
	\small
	\begin{longtable}{@{}p{0.9cm}p{5.6cm}rp{6.6cm}@{}}
		\toprule
		\textbf{Cód.} & \textbf{Criterio} & \textbf{Pts.} & \textbf{Descriptor de nivel máximo} \\
		\midrule
		\endhead
		C1 & Rigor del diseño experimental & 18 & Los tres experimentos completos, con control negativo validado, semillas fijadas y cinco repeticiones \\
		\addlinespace
		C2 & Resultado de seguridad y su análisis & 15 & Tabla de doble adjudicación con intervalos de confianza binomiales, y discusión honesta de la diferencia entre cero observado y cero garantizado \\
		\addlinespace
		C3 & Medición de la recuperación del árbitro & 10 & Tiempos de detección y reelección, solicitudes perdidas y duplicadas, y verificación de que la seguridad se mantuvo \\
		\addlinespace
		C4 & Reproducibilidad & 12 & Otro equipo regenera todas las figuras ejecutando el \emph{notebook}, sin intervención manual \\
		\addlinespace
		C5 & Cobertura de los temas de la asignatura & 14 & Los 27 temas de la tabla de la sección 7 con evidencia concreta y localizable \\
		\addlinespace
		C6 & Arquitectura distribuida y su justificación & 9 & Fronteras argumentadas, posición CAP por agregado sostenida con evidencia propia del experimento 3 \\
		\addlinespace
		C7 & Calidad del software y pruebas & 7 & Cobertura $\geq 70\,\%$ del motor de reservas, integración y carga \\
		\addlinespace
		C8 & Documento LaTeX y bibliografía & 8 & Compila sin advertencias graves ni referencias sin resolver; bibliografía verificada y anclada al dominio \\
		\addlinespace
		C9 & Borrador del manuscrito & 7 & Título e identidad nuevos aplicados, estructura completa y al menos cinco referencias de la revista objetivo \\
		\bottomrule
	\end{longtable}
\end{center}

\subsection{Penalizaciones individuales}

\begin{itemize}[leftmargin=*,itemsep=3pt]
	\item Integrante sin \emph{commits} sustantivos: se le descuenta el 50\,\% de la nota del grupo.
	\item Integrante que no responde ninguna pregunta en la defensa: se le descuenta el 30\,\%.
	\item Indicios de texto generado por IA presentado como propio sin declararlo en E19: se aplica el criterio de piso P7 y se remite el caso según la normativa de integridad académica.
\end{itemize}

\section{Lista de verificación final}

\begin{center}
	\small
	\begin{tabular}{@{}p{0.7cm}p{14.3cm}@{}}
		\toprule
		\textbf{$\square$} & \textbf{Verificación} \\
		\midrule
		$\square$ & El repositorio, el código, los contenedores y el documento usan el nombre \textsc{LabQuorum} \\
		$\square$ & La comprobación de colisiones del nombre está documentada en el \texttt{README.md} \\
		$\square$ & El PDF de carátula con la URL del repositorio en una sola línea está subido al SGA \\
		$\square$ & Clonamos el repositorio en una máquina limpia y el \texttt{.tex} compiló siguiendo el \texttt{README.md} \\
		$\square$ & \texttt{docker compose up} levanta el sistema y los veinte agentes simulados \\
		$\square$ & Cambiar \texttt{ARBITER} conmuta entre las cinco estrategias sin tocar código \\
		$\square$ & \textbf{La estrategia S0 produce dobles adjudicaciones bajo contención máxima} (control negativo validado) \\
		$\square$ & Están los datos crudos de los tres experimentos y el informe del oráculo \\
		$\square$ & La tabla de seguridad reporta intervalos de confianza binomiales, no porcentajes desnudos \\
		$\square$ & La caída del coordinador se provoca en un momento controlado y está medida \\
		$\square$ & El \emph{notebook} regenera todas las figuras sin intervención manual \\
		$\square$ & La curva de \emph{speedup} tiene los cuatro puntos y la fracción de Amdahl está calculada \\
		$\square$ & El tablero de Grafana está versionado como código \\
		$\square$ & La cobertura de pruebas del motor de reservas es $\geq 70\,\%$ \\
		$\square$ & Todos los DOI de la bibliografía resuelven y ninguno es de control industrial ni de redes eléctricas \\
		$\square$ & El borrador del manuscrito cita cinco o más artículos de la revista objetivo \\
		$\square$ & La declaración de uso de IA está completa y firmada por los cuatro integrantes \\
		$\square$ & El DOI de Zenodo resuelve y el paquete contiene generador, inyector, agentes, datos y oráculo \\
		\bottomrule
	\end{tabular}
\end{center}

\vspace{8pt}
\begin{aviso}{Una última observación para el equipo}
	Este proyecto llega a la entrega final con la mejor pregunta de investigación de la asignatura, y esa ventaja no se hereda: se pierde si el experimento se improvisa en la última semana. La diferencia entre un trabajo de aula notable y un manuscrito enviable está en tres cosas concretas que caben en dos semanas: el control negativo, las repeticiones y la medición de la reelección. Si esas tres salen bien, el resto es redacción.
\end{aviso}

\nocite{*}
\bibliographystyle{ieeetr}
\bibliography{references_PFC_LabQuorum}

\end{document}
